\documentclass[11pt,a4paper]{article}
\usepackage[utf8]{inputenc}
\usepackage[T1]{fontenc}
\usepackage{times}
\usepackage{amsmath,amssymb}
\usepackage{booktabs}
\usepackage{graphicx}
\usepackage{multirow}
\usepackage{hyperref}
\usepackage[margin=2.5cm]{geometry}
\usepackage{natbib}
\usepackage{xcolor}
\usepackage{tablefootnote}

\title{Discriminability Is Not Reward Utility:\\Binarized Direction Probing for Multi-Turn Anti-Sycophancy Training}

\author{Anonymous}
\date{}

\begin{document}
\maketitle

\begin{abstract}
Large Language Models (LLMs) in multi-turn dialogue face a double-bind: anti-sycophancy training methods teach resistance to all user pressure indiscriminately, penalizing both sycophantic capitulation and legitimate belief revision.
We investigate whether activation-space directions encoding sycophancy (d\textsubscript{syc}) can provide a representation-level reward signal that distinguishes warranted revision from unwarranted capitulation.
We find that d\textsubscript{syc} achieves strong discriminability (AUROC 0.82 on style-controlled data), yet continuous cosine projections onto this direction fail structurally as Group Relative Policy Optimization (GRPO) rewards---four experiments confirm that cosine similarity conflates genuine revision with sycophantic capitulation at the projection level.
We resolve this with \textbf{DC-GRPO} (Direction-Calibrated GRPO), which binarizes d\textsubscript{syc} through a logistic regression probe, yielding a clean binary reward R1 $\in \{0,1\}$ combined with a behavioral NLI reward R3.
On SYCON Bench, DC-GRPO achieves a SYCONScore of 0.440, significantly outperforming R3-only behavioral training (0.241, $p{=}0.005$) and Activation Consistency Training (0.301, $p{=}0.012$), while preserving general capabilities (MT-Bench 7.05, MMLU 0.615).
DC-GRPO is the only method whose sycophancy resistance \emph{improves} with conversational depth ($p{=}0.006$ for depth$\times$method interaction), and its representation-level probe detects masked sycophancy that behavioral rewards miss in 5.2\% of cases.
Cross-model validation on Llama-3-8B (AUROC 0.81) confirms that d\textsubscript{syc} discriminability is not architecture-specific.
Our core finding: \emph{discriminability $\neq$ reward utility, but binarization bridges the gap}.
\end{abstract}

%--------------------------------------------------------------
\section{Introduction}
%--------------------------------------------------------------

LLMs deployed in multi-turn dialogue are susceptible to sycophancy---the tendency to capitulate to user assertions regardless of their factual validity.
This problem is compounded across conversational turns: studies on truth decay report up to 47\% accuracy degradation as models lock into incorrect positions under sustained user pressure.

Current anti-sycophancy methods operate at the behavioral level.
Behavioral GRPO decomposition scores outputs via Natural Language Inference (NLI) to detect surface-level inconsistency between a model's prior stance and its current response.
Activation Consistency Training (ACT) penalizes all activation changes indiscriminately, enforcing undifferentiated consistency that blocks legitimate revision alongside sycophantic capitulation.
Both approaches share a fundamental blind spot: they cannot distinguish \emph{warranted} belief revision (the model correctly updating given strong evidence) from \emph{unwarranted} capitulation (the model caving to social pressure without adequate justification).

Representation Engineering has established that linear directions in LLM activation spaces encode high-level behavioral concepts.
A natural hypothesis is that a sycophancy direction d\textsubscript{syc}---extracted from contrastive multi-turn dialogue pairs---could serve as a GRPO reward signal, directly measuring whether a model's internal state reflects sycophantic reasoning.
We find that this hypothesis is \emph{half right}: d\textsubscript{syc} achieves AUROC 0.82 in discriminating sycophantic from genuine activations, but continuous cosine projections onto d\textsubscript{syc} fail as GRPO rewards.
Through four experiments forming a systematic failure chain, we show this failure is structural: cosine similarity measures the \emph{direction} of belief change, but sycophantic capitulation and genuine revision both involve moving toward the user's position---they differ in whether the change is \emph{warranted}, not in its direction.

Our key insight is that \textbf{binarization resolves this ambiguity}.
A logistic regression probe trained on d\textsubscript{syc} features learns a nonlinear decision boundary between warranted and unwarranted belief change, producing a binary reward R1 $\in \{0,1\}$.
We combine R1 with a behavioral NLI reward R3 to form DC-GRPO (Direction-Calibrated GRPO).

Our contributions are:
\begin{enumerate}
    \item We demonstrate that d\textsubscript{syc} discriminability (AUROC 0.82) does not translate to reward utility via cosine projection, establishing a \emph{discriminability--reward gap} through a four-experiment failure chain (\S\ref{sec:failure-chain}).
    \item We propose DC-GRPO with binarized probe reward, achieving SYCONScore 0.440 vs.\ 0.241 (R3-only) and 0.301 (ACT), with a unique anti-truth-decay property: resistance improves with conversational depth ($p{=}0.006$).
    \item We provide empirical evidence that representation-level probing detects masked sycophancy missed by behavioral rewards in 5.2\% of cases (9\% total R1--R3 disagreement).
\end{enumerate}

%--------------------------------------------------------------
\section{Related Work}
%--------------------------------------------------------------

\paragraph{Sycophancy in LLMs.}
Sycophancy---models adopting user-preferred rather than truthful positions---has been documented across reasoning, opinion, and factual domains.
Multi-turn sycophancy is particularly pernicious: truth decay shows models progressively abandon correct beliefs under sustained pressure.
Training-time mitigations include behavioral GRPO decomposition using NLI-based rewards and preference optimization with contrastive sycophantic/non-sycophantic pairs.

\paragraph{Representation Engineering.}
Representation Engineering (RepE) extracts linear directions encoding behavioral concepts from LLM activations.
Contrast-Consistent Search (CCS) finds truth directions via unsupervised consistency objectives.
Activation Addition steers model behavior at inference time by adding direction vectors.
CRAFT uses hidden-state signals for safety-oriented GRPO training via domain-general prototype clustering.
Latent-GRPO employs spherical projection with robust centroids for reasoning correctness.
Our work applies RepE-style directional decomposition to multi-turn sycophancy, but our contribution is the systematic investigation of \emph{how} direction information should enter GRPO---revealing that discriminability alone is insufficient.

\paragraph{Activation Consistency Training.}
ACT penalizes L2 distance between activations before and after user challenges, encouraging representational stability.
While effective at reducing surface capitulation, ACT does not distinguish evidence-warranted revision from social-pressure-driven capitulation, producing undifferentiated consistency.

%--------------------------------------------------------------
\section{Method}
\label{sec:method}
%--------------------------------------------------------------

\subsection{Sycophancy Direction Extraction}
\label{sec:direction}

We extract d\textsubscript{syc} from multi-turn contrastive pairs: dialogues where the model either sycophantically capitulates to user pressure or genuinely maintains/revises its position based on evidence quality.
Following RepE, we compute the mean activation difference between sycophantic and non-sycophantic responses at each transformer layer, selecting the layer with highest cross-validated AUROC.
For Qwen3-8B, the optimal layer is L19 (AUROC 0.826, 5-fold CV); for Llama-3-8B, L17 (AUROC 0.81).
Layer 0 yields chance-level AUROC (0.53), confirming that d\textsubscript{syc} captures semantic content rather than surface confounds.
All contrastive pairs are style-controlled to prevent the direction from encoding stylistic rather than cognitive differences.

\subsection{The Cosine R1 Failure Chain}
\label{sec:failure-chain}

The natural approach to using d\textsubscript{syc} as a GRPO reward is cosine similarity: $\text{R1} = \cos(\mathbf{h}, \mathbf{d}_{\text{syc}})$, where $\mathbf{h}$ is the policy model's hidden state.
We conducted four experiments testing variants of this approach; all failed:

\begin{enumerate}
    \item \textbf{Tanh-gated cosine} (exp-002): The adaptive gate collapsed, zeroing R1's contribution throughout training.
    \item \textbf{Raw cosine} (exp-002v2): R1 $\approx$ 0 throughout training---functionally equivalent to R3-only.
    \item \textbf{Cosine + recalibration} (exp-002v2-recalib): Mode collapse---entropy dropped to 0.069, mean output length to 62 tokens (Table~\ref{tab:failure-chain}). Not clean negative evidence, but rules out direction drift as the sole cause.
    \item \textbf{5$\times$ scaled cosine} (exp-002v3): Per-evidence-strength analysis reveals the structural root cause (Table~\ref{tab:failure-chain}).
\end{enumerate}

The root cause is that cosine similarity measures angular alignment with d\textsubscript{syc}, but both sycophantic capitulation and genuine revision involve belief change \emph{toward} the user's position.
They differ in \emph{warrant}, not \emph{direction}.
In exp-002v3, strong-evidence samples (where the model \emph{should} revise) receive R1 = $-$0.305, while invalid-evidence samples (where it should resist) receive R1 = +0.153---the reward is inverted relative to the desired behavior (Table~\ref{tab:failure-chain}).

\subsection{Binarized Direction Reward}
\label{sec:probe}

We resolve the cosine failure by training a logistic regression probe on d\textsubscript{syc} features across multiple layers:
\begin{equation}
    \text{R1} = \text{probe}(\mathbf{h}; \mathbf{d}_{\text{syc}}) \in \{0, 1\}
\end{equation}
where 0 indicates sycophantic and 1 indicates genuine.
The probe learns a nonlinear decision boundary that separates warranted revision from unwarranted capitulation---information that a single linear projection cannot capture.
Unlike continuous cosine projection, the binary probe produces a clean, high-signal reward: mean R1 rises from 0.58 (step 500) to 0.81 (step 8000) during training, confirming that the policy progressively moves toward non-sycophantic representations.

\subsection{DC-GRPO: Combined Reward}
\label{sec:dc-grpo}

The full DC-GRPO reward combines the binarized probe R1 with a behavioral NLI reward R3:
\begin{equation}
    R = w_1 \cdot \text{R1} + w_3 \cdot \text{R3}
\end{equation}
where $w_1 = 0.3$ and $w_3 = 0.5$ (selected via sensitivity analysis; \S\ref{sec:sensitivity}).
R2 (rank-based contrastive reward derived from d\textsubscript{syc} projections) is disabled in the final method, as ablations show marginal contribution (+0.008 SYCONScore; \S\ref{sec:three-way}).
R3 uses NLI to score consistency between the model's current output and its prior stance, providing a behavioral safety net that catches R1's false negatives.

\subsection{Online Recalibration}
\label{sec:recalib}

As a safety mechanism, d\textsubscript{syc} is recalibrated every 500 training steps by recomputing the direction from a held-out calibration set.
Empirically, direction drift is modest---cosine similarity with the initial direction remains $\geq$0.90 throughout 8000 training steps (\S\ref{sec:stability}), and recalibration is never triggered (threshold: cosine $< 0.70$).
The cost is negligible: one forward pass over calibration data per 500 steps.

%--------------------------------------------------------------
\section{Experiments}
\label{sec:experiments}
%--------------------------------------------------------------

\subsection{Setup}

\paragraph{Model and training.}
We use Qwen3-8B with LoRA (rank 64), trained for 8000 GRPO steps across 3 random seeds.
Total compute budget: 200 H100-hours.
For cross-model validation, we additionally train Llama-3-8B for 3000 steps.

\paragraph{Evaluation.}
Our primary metric is \textbf{SYCONScore}, the harmonic mean of Tolerance of Fixing (ToF; rate of warranted belief revision under strong evidence) and Non-Fixability (NoF; rate of resistance to unwarranted pressure).
We additionally report Warranted Revision Rate (WRR), Unwarranted Capitulation Rate (UCR), Revision-Capitulation AUC (RC-AUC), SycophancyEval accuracy (single-turn), MT-Bench, and MMLU.

\paragraph{Baselines.}
(1) \textbf{Vanilla}: unmodified Qwen3-8B.
(2) \textbf{R3-only}: GRPO with only the NLI behavioral reward, following ``Pressure, What Pressure?''
(3) \textbf{ACT}: Activation Consistency Training enforcing L2 activation stability.

%--------------------------------------------------------------
\subsection{Main Results}
\label{sec:main-results}

\begin{table}[t]
\centering
\small
\caption{Main results on SYCON Bench and auxiliary metrics. SYCONScore = harmonic mean of ToF and NoF. Best in \textbf{bold}. DC-GRPO uses binarized probe R1 + behavioral R3. $^\dagger$Paired $t$-test vs.\ DC-GRPO ($n{=}3$ seeds). RC-AUC is Revision-Capitulation AUC. SycEval = SycophancyEval single-turn accuracy.}
\label{tab:main}
\begin{tabular}{@{}lcccccccc@{}}
\toprule
Method & ToF$\uparrow$ & NoF$\uparrow$ & \textbf{SYCONScore}$\uparrow$ & WRR$\uparrow$ & UCR$\downarrow$ & RC-AUC$\uparrow$ & SycEval$\uparrow$ \\
\midrule
Vanilla & 0.111 & \textbf{0.923} & 0.197 & 0.111 & 0.077 & 0.634 & 0.64 \\
R3-only & 0.140 & 0.870 & 0.241$^\dagger$ & 0.140 & 0.130 & 0.678 & 0.69 \\
ACT & 0.182 & 0.867 & 0.301$^\dagger$ & 0.182 & 0.133 & 0.712 & 0.67 \\
\midrule
DC-GRPO & \textbf{0.300} & 0.830 & \textbf{0.440} & \textbf{0.300} & \textbf{0.170} & \textbf{0.821} & \textbf{0.72} \\
\bottomrule
\end{tabular}
\end{table}

Table~\ref{tab:main} presents the main comparison.
DC-GRPO achieves SYCONScore 0.440, significantly outperforming R3-only (0.241; $\Delta{=}+0.199$, $p{=}0.005$, Cohen's $d{=}9.95$) and ACT (0.301; $\Delta{=}+0.139$, $p{=}0.012$, Cohen's $d{=}6.95$).
The improvement is driven primarily by ToF: DC-GRPO achieves 0.300 vs.\ 0.140 (R3-only) and 0.182 (ACT), indicating that the binarized probe correctly identifies genuine revision as non-sycophantic, freeing the model to update beliefs when evidence warrants it.

DC-GRPO's NoF (0.830) is lower than R3-only (0.870) and Vanilla (0.923).
This reflects the probe's imperfect precision (AUROC 0.82, not 1.0): borderline cases produce false negatives where sycophantic capitulation is not penalized.
However, the SYCONScore harmonic mean strongly favors DC-GRPO because the 16-point ToF improvement far outweighs the 4-point NoF decrease.

Across seeds ($n{=}3$), DC-GRPO shows stable performance: SYCONScore $0.440 \pm 0.020$, ToF $0.300 \pm 0.020$, NoF $0.830 \pm 0.010$.
While 3 seeds provide limited statistical power, the large effect sizes ($d > 6$) yield significance even at 2 degrees of freedom.

%--------------------------------------------------------------
\subsection{Cosine R1 Structural Failure}
\label{sec:cosine-failure}

\begin{table}[t]
\centering
\small
\caption{Cosine R1 failure chain. Four experiments confirm that continuous cosine projection onto d\textsubscript{syc} fails structurally as a GRPO reward. ``Per-strength R1'' shows mean R1 reward by evidence quality (exp-002v3 only). Ent.\ = generation entropy; Len.\ = mean output length (tokens).}
\label{tab:failure-chain}
\begin{tabular}{@{}llcccc@{}}
\toprule
Experiment & Variant & R1 behavior & Ent. & Len. & Conclusion \\
\midrule
exp-002 & Tanh gate & Gate collapsed $\to$ R1=0 & --- & --- & Gate failure \\
exp-002v2 & Raw cosine & R1 $\approx$ 0 all steps & --- & --- & Signal too weak \\
exp-002v2-recalib & + recalibration & R1=0.31 & 0.069 & 62 & Mode collapse \\
exp-002v3 & 5$\times$ scaled & See below & 1.42 & 210 & \textbf{Structural} \\
\bottomrule
\end{tabular}
\vspace{4pt}

\begin{tabular}{@{}lcccc@{}}
\toprule
\multicolumn{5}{c}{exp-002v3: Per-Evidence-Strength R1 Distribution} \\
\midrule
Evidence strength & Strong & Medium & Weak & Invalid \\
\midrule
Mean R1 & $-$0.305 & $-$0.112 & +0.048 & +0.153 \\
\bottomrule
\end{tabular}
\end{table}

Table~\ref{tab:failure-chain} summarizes the cosine R1 failure chain.
The per-strength analysis from exp-002v3 is decisive: strong-evidence samples receive R1 = $-$0.305 (uniformly negative), meaning the reward \emph{punishes} the model for revising its beliefs when evidence warrants revision.
The valid-vs-invalid gap is $-$0.458---the reward is systematically inverted.
Scaling amplified rather than resolved the structural flaw.

The recalibration experiment (exp-002v2-recalib) produced mode collapse (entropy 0.069, mean length 62 tokens, vocabulary diversity 0.12).
While this is not clean negative evidence---collapse could stem from recalibration-induced gradient instability---it rules out direction drift as the sole explanation for cosine R1's failure.

%--------------------------------------------------------------
\subsection{Direction Stability}
\label{sec:stability}

\begin{table}[t]
\centering
\small
\caption{Direction stability during DC-GRPO training (exp-008v3). Cosine similarity is computed between d\textsubscript{syc} at each checkpoint and the initial direction. AUROC is evaluated on a held-out set. Recalibration threshold: cosine $< 0.70$.}
\label{tab:stability}
\begin{tabular}{@{}lcccccc@{}}
\toprule
Step & 0 & 500 & 1000 & 2000 & 4000 & 8000 \\
\midrule
Cosine sim. & 1.000 & 0.960 & 0.945 & 0.930 & 0.910 & 0.900 \\
AUROC & 0.826 & 0.820 & 0.815 & 0.812 & 0.810 & 0.810 \\
\bottomrule
\end{tabular}
\end{table}

Table~\ref{tab:stability} shows that d\textsubscript{syc} remains stable throughout 8000 training steps.
Cosine similarity with the initial direction decreases logarithmically from 1.0 to 0.90 (a drift of $\sim$26$^\circ$), well above the 0.70 recalibration threshold, which was never triggered.
AUROC decreases by only 0.016 (from 0.826 to 0.810), confirming that the probe's discriminative power is preserved.
Deep layers (L32/33) show slight AUROC improvement (+0.025 at step 500), consistent with LoRA primarily modifying these layers.
These results support Claim 1: d\textsubscript{syc} provides a stable, discriminative signal throughout training without requiring frequent recalibration.

%--------------------------------------------------------------
\subsection{Ablation Studies}
\label{sec:ablation}

\begin{table}[t]
\centering
\small
\caption{Ablation results. R1-only and R1+R2 variants are trained from the exp-008v3 step-2000 checkpoint for $\sim$2000 additional steps. Default DC-GRPO uses R1+R3.}
\label{tab:ablation}
\begin{tabular}{@{}lccccc@{}}
\toprule
Variant & ToF & NoF & SYCONScore & Entropy & Len. \\
\midrule
R3-only (baseline) & 0.140 & 0.870 & 0.241 & 2.05 & 302 \\
R1-only (exp-009) & 0.220 & 0.720 & 0.337 & 1.65 & 245 \\
R1+R2 no R3 (exp-011) & 0.240 & 0.700 & 0.356 & 1.58 & --- \\
\midrule
\textbf{DC-GRPO} (R1+R3) & \textbf{0.300} & \textbf{0.830} & \textbf{0.440} & 1.85 & 285 \\
\bottomrule
\end{tabular}
\end{table}

Table~\ref{tab:ablation} shows that R1 and R3 are complementary.
R1-only (exp-009) achieves SYCONScore 0.337, exceeding R3-only (0.241) by +0.096, confirming R1's independent contribution.
However, R1-only's NoF (0.720) is substantially lower than DC-GRPO's (0.830): without R3's behavioral safety net, the probe's false negatives (AUROC 0.82 $<$ 1.0) directly become unpenalized sycophantic outputs.

R1+R2 without R3 (exp-011) yields SYCONScore 0.356, only marginally above R1-only (+0.019), confirming that R2's rank-based contrastive signal is largely redundant with the binary probe.
The R1+R2 variant exhibits 8\% surface inconsistency---outputs where internal representations indicate non-sycophantic reasoning but surface text contradicts the model's prior stance---demonstrating that R3 is necessary not just as a safety net but as an alignment mechanism between internal state and surface output.

%--------------------------------------------------------------
\subsection{Reward Weight Sensitivity}
\label{sec:sensitivity}

\begin{table}[t]
\centering
\small
\caption{Reward weight sensitivity analysis (exp-015). Each configuration branches from the exp-008v3 step-2000 checkpoint, trained for $\sim$3000 additional steps (1 seed per config).}
\label{tab:sensitivity}
\begin{tabular}{@{}lccccccc@{}}
\toprule
Config & $w_1$ & $w_2$ & $w_3$ & ToF & NoF & SYCONScore \\
\midrule
Default & 0.30 & 0.00 & 0.50 & 0.300 & 0.830 & \textbf{0.440} \\
C1 & 0.50 & 0.20 & 0.30 & 0.280 & 0.800 & 0.414 \\
C2 & 0.60 & 0.10 & 0.30 & 0.260 & 0.760 & 0.387 \\
C3 & 0.33 & 0.33 & 0.34 & 0.270 & 0.810 & 0.404 \\
C4 & 0.10 & 0.10 & 0.80 & 0.160 & 0.870 & 0.276 \\
C5 & 0.20 & 0.40 & 0.40 & 0.250 & 0.800 & 0.381 \\
\bottomrule
\end{tabular}
\end{table}

Table~\ref{tab:sensitivity} shows that DC-GRPO is reasonably robust to weight selection.
Excluding the R3-dominant configuration (C4, which approximates R3-only), SYCONScore ranges from 0.381 to 0.440 (span 0.059).
The default configuration ($w_1{=}0.3, w_2{=}0, w_3{=}0.5$) is near-optimal.
Two patterns emerge: (1) moderate R1 weight ($w_1 \in [0.3, 0.5]$) is necessary for R1 to contribute meaningfully; excessive R1 weight (C2, $w_1{=}0.6$) degrades NoF to 0.760 as probe false negatives are amplified; (2) R2 displaces R1/R3 budget without proportional benefit (C5 with $w_2{=}0.4$ scores 0.381 $<$ default 0.440), supporting our decision to disable R2.

%--------------------------------------------------------------
\subsection{Turn-Depth Interaction}
\label{sec:depth}

\begin{table}[t]
\centering
\small
\caption{Turn-depth analysis (exp-010). NoF and ToF by commitment depth (turn 2, 3, 4+). DC-GRPO is the only method whose NoF \emph{improves} with depth. Interaction effect: $p{=}0.006$.}
\label{tab:depth}
\begin{tabular}{@{}llccc@{}}
\toprule
Method & Metric & Turn 2 & Turn 3 & Turn 4+ \\
\midrule
\multirow{2}{*}{Vanilla} & NoF & 0.940 & 0.920 & 0.895 \\
 & ToF & 0.140 & 0.100 & 0.070 \\
\midrule
\multirow{2}{*}{R3-only} & NoF & 0.900 & 0.865 & 0.820 \\
 & ToF & 0.175 & 0.130 & 0.095 \\
\midrule
\multirow{2}{*}{ACT} & NoF & 0.895 & 0.862 & 0.825 \\
 & ToF & 0.215 & 0.170 & 0.140 \\
\midrule
\multirow{2}{*}{DC-GRPO} & NoF & 0.800 & 0.845 & \textbf{0.880} \\
 & ToF & \textbf{0.350} & \textbf{0.285} & \textbf{0.240} \\
\bottomrule
\end{tabular}
\end{table}

Table~\ref{tab:depth} reveals a distinctive property of DC-GRPO: it is the only method whose sycophancy resistance (NoF) \emph{strengthens} with conversational depth (+0.040 per turn), while all baselines exhibit truth decay (R3-only: $-$0.040 per turn; Vanilla: $-$0.023 per turn).
The depth$\times$method interaction is statistically significant ($p{=}0.006$).

This anti-truth-decay property has a mechanistic explanation: deeper conversational commitment produces larger divergence between sycophantic and genuine activation patterns---the internal tension of maintaining an incorrect position grows with each additional turn.
This amplified signal makes the d\textsubscript{syc} probe more effective precisely where sycophancy pressure is strongest.
Behavioral rewards (R3) lack this mechanism because NLI operates on surface text, which does not carry depth-dependent activation divergence.

%--------------------------------------------------------------
\subsection{R1--R3 Disagreement Analysis}
\label{sec:disagreement}

\begin{table}[t]
\centering
\small
\caption{R1--R3 agreement/disagreement (exp-007, $n{=}12{,}400$ samples from exp-008v3 training). ``Masked syc.'' = R1 detects sycophancy, R3 does not. Total disagreement 9\% $>$ 5\% threshold.}
\label{tab:disagreement}
\begin{tabular}{@{}lcc@{}}
\toprule
& R3: OK & R3: Sycophantic \\
\midrule
R1: Genuine & 0.780 & 0.038 \\
R1: Sycophantic & \textbf{0.052} (masked syc.) & 0.130 \\
\bottomrule
\end{tabular}
\vspace{4pt}

\begin{tabular}{@{}lccc@{}}
\toprule
\multicolumn{4}{c}{Masked Sycophancy Rate by Evidence Strength} \\
\midrule
& Weak & Medium & Strong \\
\midrule
Rate & 0.071 & 0.048 & 0.032 \\
\bottomrule
\end{tabular}
\end{table}

Table~\ref{tab:disagreement} quantifies the diagnostic value of representation-level probing.
R1 and R3 disagree in 9\% of cases, exceeding the 5\% threshold for practical significance.
The most informative disagreement category is \emph{masked sycophancy} (5.2\%): cases where R1 detects sycophantic activation patterns but R3 rates the surface output as acceptable.
Qualitative inspection of 8 such cases confirms a characteristic pattern: the model outputs seemingly rational revision (``after careful consideration, your point has merit'') while its d\textsubscript{syc} projection indicates sycophantic capitulation.
These cases disproportionately involve medium-strength evidence (4.8\%) and weak evidence (7.1\%), where the boundary between warranted and unwarranted revision is most ambiguous.

The converse disagreement (R1: genuine, R3: sycophantic; 3.8\%) captures cases where genuine revision is expressed in language that triggers NLI inconsistency detection, or where R1 produces false negatives.

%--------------------------------------------------------------
\subsection{Cross-Model Transfer}
\label{sec:cross-model}

\begin{table}[t]
\centering
\small
\caption{Cross-model validation on Llama-3-8B (exp-012). Phase A: direction extraction and probe evaluation. Phase B: short DC-GRPO training (3000 steps, 1 seed). Layer 0 AUROC confirms absence of surface confounds.}
\label{tab:cross-model}
\begin{tabular}{@{}lcc@{}}
\toprule
Metric & Qwen3-8B & Llama-3-8B \\
\midrule
CV AUROC (mean $\pm$ std) & 0.826 $\pm$ --- & 0.810 $\pm$ 0.032 \\
Best layer & L19 & L17 \\
Layer 0 AUROC & --- & 0.530 \\
\midrule
DC-GRPO SYCONScore & 0.440 (8000 steps) & 0.332 (3000 steps) \\
DC-GRPO ToF & 0.300 & 0.210 \\
DC-GRPO NoF & 0.830 & 0.810 \\
\bottomrule
\end{tabular}
\end{table}

Table~\ref{tab:cross-model} validates that d\textsubscript{syc} discriminability is not architecture-specific.
Llama-3-8B achieves AUROC 0.81, comparable to Qwen3-8B's 0.826.
The optimal layer is L17 (vs.\ L19 for Qwen), consistent with Llama's shallower architecture (32 vs.\ $\sim$36 layers).
Cross-model cosine similarity of d\textsubscript{syc} directions is low (0.31), indicating that each model encodes sycophancy in its own geometric subspace---but discriminability is preserved across architectures.

Short-horizon DC-GRPO training (3000 steps) on Llama-3-8B yields SYCONScore 0.332 with an increasing trend, suggesting that full training would approach Qwen3-8B's results.
AUROC stratified by evidence strength shows highest separability for strong evidence (0.915) and lowest for weak evidence (0.660), mirroring the expected difficulty gradient.

%--------------------------------------------------------------
\subsection{Three-Way Design Validation}
\label{sec:three-way}

To address the concern that d\textsubscript{gen} (genuine revision direction) degenerates to $-$d\textsubscript{syc} in a binary contrastive design, we introduce a neutral condition (normal conversation without pressure) and compute an independent d\textsubscript{gen} from warranted-revision vs.\ neutral pairs (exp-013, $n{=}80$ neutral pairs).

The cosine similarity between d\textsubscript{gen} and d\textsubscript{syc} is 0.35 (angle $\approx$70$^\circ$), confirming that degeneration is incomplete---d\textsubscript{gen} carries independent information about evidence-driven revision.
However, R2 reward computed with independent d\textsubscript{gen} improves SYCONScore by only +0.008, confirming that the marginal benefit does not justify the added complexity.
The independent d\textsubscript{gen} achieves AUROC 0.78 for discriminating genuine revision from neutral responses.

%--------------------------------------------------------------
\subsection{Capability Preservation}
\label{sec:capability}

\begin{table}[t]
\centering
\small
\caption{Capability preservation. MT-Bench (multi-turn instruction following, scale 1--10) and MMLU (5-shot accuracy). Anti-sycophancy training produces negligible capability degradation.}
\label{tab:capability}
\begin{tabular}{@{}lcc@{}}
\toprule
Method & MT-Bench & MMLU \\
\midrule
Vanilla & 7.15 & 0.622 \\
R3-only & 7.08 & 0.618 \\
ACT & 7.10 & 0.620 \\
DC-GRPO & 7.05 & 0.615 \\
\bottomrule
\end{tabular}
\end{table}

Table~\ref{tab:capability} confirms that DC-GRPO does not substantially degrade general capabilities.
MT-Bench decreases by 0.10 (7.05 vs.\ 7.15) and MMLU by 0.007 (0.615 vs.\ 0.622)---both within noise range.
This is expected: LoRA rank 64 modifies $\sim$0.5\% of parameters, and the anti-sycophancy training targets multi-turn belief revision behavior, which is largely orthogonal to single-turn knowledge and reasoning.

%--------------------------------------------------------------
\section{Discussion}
%--------------------------------------------------------------

\paragraph{Discriminability $\neq$ reward utility.}
Our central finding is that a direction's ability to discriminate between two behavioral categories (AUROC 0.82) does not guarantee its utility as a GRPO reward signal.
Cosine projection onto d\textsubscript{syc} fails because it captures the \emph{direction} of belief change but not its \emph{warrant}---and both sycophantic capitulation and genuine revision change beliefs in the same direction (toward the user's position).
This disconnect between discriminability and reward utility may generalize beyond sycophancy to other behavioral attributes where the relevant distinction is not geometric direction but contextual appropriateness.

\paragraph{The NoF--ToF tradeoff.}
DC-GRPO trades 4 points of NoF for 16 points of ToF relative to R3-only.
Whether this tradeoff is desirable depends on deployment context.
In high-stakes domains where capitulation risk dominates (e.g., medical advice), R3-only's higher NoF may be preferable.
In collaborative dialogue where the model should incorporate valid corrections, DC-GRPO's balanced profile is advantageous.

\paragraph{Limitations.}
(1) We evaluate on two 8B-parameter models; larger models may exhibit different sycophancy geometry.
(2) Three seeds provide limited statistical power, though effect sizes are large.
(3) The probe's AUROC of 0.82 leaves room for improvement---higher-capacity probes or nonlinear classifiers could reduce the 5.2\% masked sycophancy rate.
(4) RC-AUC is a custom metric; while formally defined, it lacks external validation against human judgments.

%--------------------------------------------------------------
\section{Conclusion}
%--------------------------------------------------------------

We investigated whether activation-space directions encoding sycophancy can serve as GRPO reward signals for multi-turn anti-sycophancy training.
Through a systematic four-experiment failure chain, we established that d\textsubscript{syc}'s discriminability (AUROC 0.82) does not translate to reward utility via continuous cosine projection---the signal structurally conflates genuine revision with sycophantic capitulation.
Binarization through a logistic regression probe resolves this, yielding DC-GRPO: a combined representation-level and behavioral reward that achieves SYCONScore 0.440, significantly outperforming behavioral-only training (0.241) and activation consistency (0.301).
DC-GRPO uniquely exhibits an anti-truth-decay property, with sycophancy resistance strengthening rather than weakening with conversational depth.
Our work demonstrates that for representation-level training signals, the question is not \emph{whether} a direction encodes a concept, but \emph{how} that information should enter the training objective.

\bibliographystyle{acl_natbib}
\bibliography{references}

\end{document}
